\par\Needspace{6\baselineskip}\medskip
\hypertarget{native-141-article101}{}\pdfbookmark[2]{101.}{article-101}
\begin{center}101.\end{center}
Führt man die Größe $\zeta$ ein und setzt $\BareRoot{(l-z)}=\dfrac my$ oder $\BareRoot{(L+z)}=\dfrac MY$ sowie
\[
\mathrm{15)}\qquad\frac{mm}{\tfrac56+l+\zeta}=h
\]
\noindent oder
\[
\mathrm{15^*)}\qquad\frac{MM}{L-\tfrac56-\zeta}=H,
\]
\noindent so nehmen die Gleichungen 13 und $13^*$ die Form an
\[
\begin{aligned}
\mathrm{16)}\qquad&\frac{(y-1)yy}{y+\tfrac19}=h\\[.6em]
\mathrm{16^*)}\qquad&\frac{(Y+1)YY}{Y-\tfrac19}=H,
\end{aligned}
\]
\noindent und sind somit genau dieselben wie die für die Ellipse erhaltenen Gleichungen (15 und $15^*$, Art. 91). Soweit also $h$ oder $H$ als bekannt gelten kann, lässt sich $y$ oder $Y$ daraus ableiten, und man erhält dann
\[
\begin{aligned}
\mathrm{17)}\qquad&z=l-\frac{mm}{yy}\\[.6em]
\mathrm{17^*)}\qquad&z=\frac{MM}{YY}-L.
\end{aligned}
\]
\noindent Daraus folgt, dass alle oben für die Ellipse angegebenen Rechenschritte ebenso für die Hyperbel gelten, bis aus dem Näherungswert von $h$ oder $H$ die Größe $y$ oder $Y$ gewonnen ist. Dann muss die Größe $\dfrac{mm}{yy}-l$ oder $L-\dfrac{MM}{YY}$, die bei der Ellipse positiv und bei der Parabel $=0$ sein musste, bei der Hyperbel negativ sein. Dieses Kriterium bestimmt somit die Art des Kegelschnitts. Für das gefundene $z$ liefert unsere Tafel den Wert von $\zeta$, aus dem sich ein berichtigter Wert von $h$ oder $H$ ergibt. Mit diesem Wert ist die Rechnung zu wiederholen, bis alles genau übereinstimmt.
\par
Sobald der wahre Wert von $z$ gefunden ist, ließe sich daraus $c$ mit der Formel $c=1+2z+2\BareRoot{(z+zz)}$ ableiten. Auch im Hinblick auf die folgenden Anwendungen empfiehlt es sich, den Hilfswinkel $n$ einzuführen, der durch die Gleichung $\operatorname{tang}2n=2\BareRoot{(z+zz)}$ zu bestimmen ist. Damit ergibt sich $c=\operatorname{tang}2n+\BareRoot{(1+\operatorname{tang}2n^2)}=\operatorname{tang}(45^0+n)$.
